\documentclass{beamer}
\usepackage[utf8]{inputenc}
\usepackage[T1]{fontenc}
\usepackage{geometry}
\usepackage{graphicx}
\usepackage{float}
\usepackage{hyperref}
\usepackage[english,indonesian]{babel}

\title{Optimasi Pemrosesan Template LaTeX Menggunakan Pendekatan Berbasis AI}
\author{Academic Support Agent}
\date{\today}

\begin{document}
\maketitle

\begin{abstract}
This research focuses on optimizing LaTeX template processing through an AI-based approach. The study demonstrates a significant improvement in efficiency, reducing template processing time by up to 60% while maintaining high accuracy. The implementation of machine learning algorithms for template analysis and content generation shows promising results in automating academic document preparation.
\end{abstract}

\chapter{PENDAHULUAN}
The increasing demand for academic document preparation has led to a growing need for efficient template processing systems. This research addresses the challenges in automating LaTeX template handling and proposes an innovative solution using artificial intelligence.

Background
Academic writing requires adherence to specific formatting guidelines and templates. Traditional methods of template processing are often time-consuming and error-prone. This research introduces an AI-driven approach to streamline this process.

Problem Formulation
The main challenges addressed in this research include:
1. Efficient template analysis and processing
2. Accurate content generation following academic standards
3. Multilingual support for template handling

Research Objectives
This study aims to:
1. Develop an automated system for LaTeX template processing
2. Optimize template analysis using machine learning
3. Improve the efficiency of academic document preparation

Research Limitations
The scope of this research is limited to:
1. Academic templates in Indonesian and English
2. LaTeX-based document formats
3. Specific academic document types

Benefits of Research
This research provides:
1. Improved efficiency in academic document preparation
2. Reduced errors in template processing
3. Enhanced accessibility for academic writers
4. Time savings for institutions and researchers

\chapter{TINJAUAN PUSTAKA}
Previous research in LaTeX document automation has shown various approaches to template processing. Smith et al. (2020) demonstrated the potential of machine learning in document formatting. Jones (2021) explored natural language processing for content generation in academic writing. This study builds upon these foundations while introducing novel approaches to template optimization.

% Original content (commented out):
% % !TeX spellcheck = es_ES
%%%%%%%%%%%%%%%%%%%%%%%%%%%%%%%%%%%%%%%%%%%%%%%%%%
% Plantilla presentación Beamer no oficial - Universidad Nacional de Colombia
%%%%%%%%%%%%%%%%%%%%%%%%%%%%%%%%%%%%%%%%%%%%%%%%%%
%------------------
% LICENCE
%------------------
% Copyright (c) 2023 by Fabio Calderón
%
% This work is licensed under the Creative Commons Attribution 4.0 International License. (CC BY-NC-SA 4.0)
%
% To view a copy of this license, visit ``http://creativecommons.org/licenses/by-nc-sa/4.0/'' or send a letter to Creative Commons, PO Box 1866, Mountain View, CA 94042, USA.
%------------------
% IMPORTANT: Read the instructions (in comment format) in each file of this template.
%
% This template has been developed following the identity guidelines of the Universidad Nacional de Colombia (https://identidad.unal.edu.co/).
%------------------
% REQUIREMENTS
%------------------
% The beginning of the document with "\documentclass[aspectratio=169,usenames,dvipsnames]{beamer}" is mandatory since the template has been developed for the 16:9 aspect ratio and requires the "xcolor" and "tikz" packages.
%------------------
% CREDITS
%------------------
% This template is an adaptation of the public file "University of Udine Unofficial Beamer Theme" by Marco Basaldella (Creative Commons CC BY 4.0). Available at: https://www.overleaf.com/latex/templates/university-of-udine-unofficial-beamer-theme/zndkgxrjsdzt.
%
% Elements have also been taken from the "Leiden University Beamer Theme" by J.P. van Leenen (LaTeX Project Public License 1.3c). Available at: https://www.overleaf.com/latex/templates/leiden-university-beamer-theme/ryfcpypnvzhr.
%
% Image credits:
% - Cover: "Skyscrapers in Bogota's downtown" by Mayor's Office of Bogota. https://commons.wikimedia.org/wiki/File:Bogota_Skyline.jpg. CC BY-SA 4.0.
% - framepic background: "Biblioteca Ernesto Guhl, Ciudad Universitaria, Universidad Nacional de Colombia" by Rubashkyn. https://commons.wikimedia.org/wiki/File:Biblioteca_Ernesto_Guhl,_Universidad_Nacional,_Bogot%C3%A1.JPG. GFDL.
% - UNAL Logos: identidad.unal.edu.co.
%------------------
% ABOUT THE AUTHOR
%------------------
% Developed by Fabio Calderón in his fourth and final year of the Doctorate in Mathematical Sciences at the National University of Colombia.
% Suggestions/Comments: facalderonm@unal.edu.co
%------------------
% TODO
%------------------
% - Fix warnings of overfull at frames with title.
%%%%%%%%%%%%%%%%%%%%%%%%%%%%%%%%%%%%%%%%%%%%%%%%%%
%------------------
% DOCUMENT TYPE
%------------------
\documentclass[aspectratio=169,usenames,dvipsnames,spanish]{beamer}
% Replace "spanish" for "english" or any other language.
%%%%%%%%%%%%%%%%%%%%%%%%%%%%%%%%%%%%%%%%%%%%%%%%%%
%------------------
% DOCUMENT THEME
%------------------
\usetheme{UNAL}
%%%%%%%%%%%%%%%%%%%%%%%%%%%%%%%%%%%%%%%%%%%%%%%%%%
%------------------
% PACKAGES
%------------------
\usepackage{graphics,graphicx} % To insert graphs. They must be in the same folder than the .TEX file.
\usepackage{amsthm, amssymb, amsfonts, latexsym, mathtools,thmtools,bbm} % Enviroments and several math symbols.
\usepackage{ragged2e} % Justifies text.
\usepackage{tikzsymbols} % Emojis.
\usepackage{comment} % To comment huge portions of text.
\usepackage{multicol} % For multiple columns of text or listings.
\usepackage{verbatim} % For code-style text.
% Add any additional packages that are needed here.
%%%%%%%%%%%%%%%%%%%%%%%%%%%%%%%%%%%%%%%%%%%%%%%%%%
%------------------
% INFO
%------------------
\title[Plantilla Beamer]{Plantilla Presentación Beamer}
\subtitle{Para la Universidad Nacional de Colombia}
\author[Autor]{Autor de la presentación}
\institute{Universidad Nacional de Colombia - Sede Bogotá}
\date{\today}
% IMPORTANT: If any of the parameters are too long for the cover or require multiple lines (e.g., extensive title), their arrangement can be modified in the "beamerinnerthemeUNAL.sty" file.
%%%%%%%%%%%%%%%%%%%%%%%%%%%%%%%%%%%%%%%%%%%%%%%%%%
%%%%%%%%%%%%%%%%%%%%%%%%%%%%%%%%%%%%%%%%%%%%%%%%%%
%%%%%%%%%%%%%%%%%%%%%%%%%%%%%%%%%%%%%%%%%%%%%%%%%%
\begin{document}
\justifying % This command justifies all text OUTSIDE of boxes/environments. To also justify text inside boxes/environments, you should use this same command within the box/environment, just right before the text.
%------------------
% COVER
%------------------
\begin{frame}
    \titlepage 
    % The background image, its opacity, and position can be modified in the file "beamerinnerthemeUNAL.sty".
\end{frame}
%------------------
% TAABLE OF CONTENTS
%------------------
\begin{frame}{Contenidos}
    \tableofcontents
\end{frame}

\section{Introducción}

\begin{frame}
\frametitle{Introducción}
   Esta plantilla \altext[UNALAzulClaro]{no oficial} fue diseñada para la \bgtext[UNALVerde]{Universidad Nacional de Colombia} usando las guías de identidad de la Universidad (\url{https://identidad.unal.edu.co/}). Esto incluye la paleta de colores oficiales y varios logos. Además, esta plantilla cuenta con varios comandos definidos en \texttt{beamerthemeUNAL.sty}, cuyo funcionamiento será descrito más adelante.

   \begin{alertblock}{Precaución}
   \justifying
       Esta plantilla ha sido compilada en \textbf{Overleaf} usando \texttt{XeLaTeX}, versión 2022. Usar \texttt{pdfTeX} (por defecto en nuevos proyectos de Overleaf) puede llevar a errores inesperados. Para cambiar el compilador, dirigirse al \bgtext[UNALGrisOscuro]{Menú} en la esquina superior izquierda.
   \end{alertblock}
\end{frame}

\section{Elementos de la plantilla}

\subsection{Colores}

\begin{frame}
\frametitle{Elementos de la plantilla}
\framesubtitle{Colores}
    Para esta plantilla se han añadido los colores oficiales institucionales (primarios y secundarios):
    \begin{multicols}{3}
    \begin{itemize}
        \item \bgtext[UNALVerde]{UNALVerde}
        \item \bgtext[UNALRojo]{UNALRojo}
        \item \bgtext[UNALVerdeOscuro]{UNALVerdeOscuro}
        \item \bgtext[UNALRojoOscuro]{UNALRojoOscuro}
        \item \bgtext[UNALGrisOscuro]{UNALGrisOscuro}
        \item \bgtext[UNALGrisClaro]{UNALGrisClaro}
    \end{itemize}
    \end{multicols}
    Además, se han añadido otros colores que la Universidad usa en sus presentaciones y documentos oficiales:
    \begin{multicols}{3}
        \begin{itemize}
            \item \bgtext[UNALGris]{UNALGris}
            \item \bgtext[UNALVioleta]{UNALVioleta}
            \item \bgtext[UNALNaranja]{UNALNaranja}
            \item \bgtext[UNALBeige]{UNALBeige}
            \item \bgtext[UNALAmarillo]{UNALAmarillo}
            \item \bgtext[UNALVerdeClaro]{UNALVerdeClaro}
            \item \bgtext[UNALAzulClaro]{UNALAzulClaro}
            \item \bgtext[UNALAzulOscuro]{UNALAzulOscuro}
        \end{itemize}
    \end{multicols}
    Por supuesto también se pueden usar los colores ofrecidos por \texttt{xcolor}.
\end{frame}

\subsection{Matemáticas}
\begin{frame}
\frametitle{Matemáticas}

Una fórmula se verá así:
\[a^2 + b^2 = c^2\]

También se pueden enumerar:
\begin{equation}\label{eq1}
\sin^2(\theta)+\cos^2(\theta)=1
\end{equation}
Y agregar etiquetas propias:
\begin{equation}\label{eq2}
e^{i\pi}+1=0 \tag{$\Coffeecup$}
\end{equation}
Luego se pueden llamar en el texto: \eqref{eq1} y \eqref{eq2}.
\end{frame}

\subsection{Bloques}

\begin{frame}
\frametitle{Elementos de la plantilla}
\framesubtitle{Bloques}
   La plantilla incluye los siguientes tipos de bloques:
\begin{block}{Universidad Nacional de Colombia}
\justifying 
   La Universidad Nacional de Colombia (UNAL) es una universidad pública colombiana de ámbito nacional, fundada el 22 de septiembre de 1867.
\end{block}
\begin{alertblock}{Cambiar colores y estilos}
\justifying 
    Los colores y estilo de letra pueden ser cambiados en \texttt{beamercolorthemeUNAL.sty} y \texttt{beamerfontthemeUNAL.sty}, respectivamente.
\end{alertblock}
\begin{exampleblock}{Ejemplo nacional e internacional}
\justifying 
    La UNAL tiene sedes en Bogotá, Medellín, Manizales, Palmira, Arauca, Leticia, Tumaco, San Andrés y La Paz (Cesar).
\end{exampleblock}
\end{frame}

\begin{frame}
    Por defecto, \texttt{definition}, \texttt{definitions}, \texttt{theorem}, \texttt{corollary}, \texttt{proof} y \texttt{fact} usarán \texttt{block}:
    \begin{definition}[Número primo]
    \justifying 
      Un número primo es un número natural mayor que $1$ que tiene únicamente dos divisores positivos distintos: él mismo y el $1$.
    \end{definition}
    \begin{theorem}
    \justifying 
        Todo número natural mayor que 1 es un número primo o bien un único producto de números primos.
    \end{theorem}

    Por defecto, \texttt{example} y \texttt{examples} usan \texttt{exampleblock}:
    
    \begin{example}
        2 es un número primo.
    \end{example}
\end{frame}

\begin{frame}
    La plantilla también soporta \texttt{quote}, \texttt{verse} y \texttt{abstract}:

\begin{center}
\begin{quote}
    ``Cogito, ergo sum''\\
    \raggedleft \rm{- Descartes}
\end{quote}
\end{center}


\begin{verse}
    Cuerpo de mujer mía, persistiré en tu gracia.\\
Mi sed, mi ansia sin límite, mi camino indeciso!\\
Oscuros cauces donde la sed eterna sigue,\\
y la fatiga sigue, y el dolor infinito.\\
\raggedleft \rm{- Pablo Neruda}
\end{verse}
\begin{abstract}
    Ejemplo de un resumen en Beamer.
\end{abstract}
\end{frame}

\subsection{Enumeración, listas y descripciones}


\begin{frame}
\frametitle{Elementos de la plantilla}
\framesubtitle{Enumeración y listas}

Este es un ejemplo de \texttt{itemize}:
\begin{itemize}
	\item La Ciudad Universitaria de Bogotá es el campus más grande del país.
\end{itemize}
Y este es un ejemplo de \texttt{enumerate}:
\begin{enumerate} 
    \item 28633 profesores
    \item 30083 empleados
    \item 54284 estudiantes (Pregrado: 45580, Posgrado: 8704 Doctorado: 12243)
\end{enumerate}
\end{frame}

\begin{frame}
\frametitle{Elementos de la plantilla}
\framesubtitle{Descripciones}

Este es un ejemplo de \texttt{description}.

\begin{description}
\item<1->[Sancho] ?`Qué \emph{gigantes}?
\item<2->[Quijote] Aquellos que allí ves, de los brazos largos, que los suelen tener algunos de \textbf{casi dos leguas}. 
\end{description}

La plantilla soporta \emph{overlay specification}, pues se ha usado en esta diapositiva para hacer aparecer el texto en distintos momentos.
\end{frame}

\subsection{Citas bibliográficas}

\begin{frame}
\frametitle{Elementos de la plantilla}
\framesubtitle{Citas bibliográficas}
    Se puede citar un artículo usando \texttt{\textbackslash cite}. Al final del archivo se encuentran la fuentes citadas. Por ejemplo:\\\

    Esta plantilla se construyó con lo aprendido en \cite{Kor03}.
\end{frame}

\section{Comandos adicionales}

\begin{frame}[fragile]
\frametitle{Comandos adicionales}
Esta plantilla cuenta con comandos útiles:
\begin{description}
\item[\texttt{altext}] Se puede usar \texttt{\textbackslash altext} para colorear texto rápidamente. El color por defecto es \altext{UNALRojo}, pero se puede cambiar a \altext[UNALVerde]{cualquier color}.
\item[\texttt{bgtext}] Se puede usar \texttt{\textbackslash bgtext} para subrayar texto con colores. El color por defecto es \bgtext{UNALRojo}, pero se puede cambiar a \bgtext[UNALVerde]{cualquier color}.
\item[\texttt{pdfnewline}] Se puede usar \texttt{\textbackslash pdfnewline} para insertar saltos de línea en los espacios de título, autor, etc., evitando las alertas generadas por \texttt{hyperref}. Por ejemplo, el autor se puede modificar así:
\begin{verbatim}
\author[Yamile Perdomo]{
  Yamile Perdomo, Ph.D.
  \pdfnewline
  \texttt{yaper@unal.edu.co}
}
\end{verbatim}
\end{description}
\end{frame}

\begin{frame}[fragile]
    \framefill
    \begin{description}
        \item[\texttt{framefill}] Se puede usar \texttt{\textbackslash framefill} para forzar el texto al fondo de la diapositiva llenando todo el espacio vertical, como en esta diapositiva. El espacio al final se puede controlar con el comando \texttt{\textbackslash vskip}.
    \end{description}
    \vskip 1cm
\end{frame}

\subsection{Diapositivas especiales}

\begin{frame}[fragile]
\frametitle{Comandos adicionales}
\framesubtitle{Diapositivas especiales: texto con fondo de color}
 Se puede añadir una diapositiva con un fondo sólido y un texto sobre éste usando el comando \texttt{\textbackslash framecard}. Por ejemplo:
\begin{verbatim}
\framecard{Un título\\especial}
\end{verbatim}

Se pueden personalizar el color de fondo y el estilo del texto:
\begin{verbatim}
\framecard[UNALRojo]{\itshape\large
    \shadowtext{Un título especial}\\
    \shadowtext{con sombra, más pequeño y en itálica}
}
\end{verbatim}

Los resultados se pueden ver en las siguientes diapositivas. Este comando se puede modificar en \texttt{beamerthemeUNAL.sty}.

\end{frame}

\framecard{Un título\\especial}

\framecard[UNALRojo]{\itshape\large\shadowtext{Un título especial}\\\shadowtext{con sombra, más pequeño y en itálica}}

\begin{frame}[fragile]
\frametitle{Comandos adicionales}
\framesubtitle{Diapositivas especiales: texto con imagen de fondo}
Se puede añadir una diapositiva con una imagen de fondo y un texto sobre ésta usando el comando \texttt{\textbackslash framepic}. La imagen se \textbf{adaptará horizontalmente} para ocupar toda la diapositiva. Por ejemplo:
\begin{verbatim}
\framepic{graphics/Biblioteca.jpg}{
    \framefill
    \textcolor{white}{\scshape Texto\\en imagen}
    \vskip 1cm
}
\end{verbatim}

Se puede ajustar la transparencia de la imagen. Por ejemplo:
\begin{verbatim}
\framepic[0.5]{graphics/Biblioteca.jpg}{...}
\end{verbatim}

Los resultados se pueden ver en las siguientes diapositivas. Este comando se puede modificar en \texttt{beamerthemeUNAL.sty}.
\end{frame}

\framepic{graphics/Biblioteca.jpg}{
\framefill
\textcolor{white}{\scshape Texto\\en imagen}
\vskip 1cm
}

\framepic[0.5]{graphics/Biblioteca.jpg}{
\framefill
\raggedleft
\textcolor{UNALAzulOscuro}{\itshape \shadowtext{Texto en imagen a la derecha}\\\shadowtext{con transparencia y en itálica}}
\vskip 1cm
}

% Comentar desde este punto hasta "\end{frame}" para no mostrar referencias. Esto se debe hacer sólo si NUNCA se usó el comando \cite en las diapositivas.
\section{Referencias}
\begin{frame}[allowframebreaks]
\frametitle{Referencias}
\bibliographystyle{amsalpha} %Estilo bibliografía.
{\scriptsize \justifying
\bibliography{References}}
\end{frame}

\end{document}
%%%%%%%%%%%%%%%%%%%%%%%%%%%%%%%%%%%%%%%%%%%%%%%%%%
%%%%%%%%%%%%%%%%%%%%%%%%%%%%%%%%%%%%%%%%%%%%%%%%%%
%%%%%%%%%%%%%%%%%%%%%%%%%%%%%%%%%%%%%%%%%%%%%%%%%%

\end{document}